\documentclass[journal]{IEEEtran}
\usepackage{amsmath,amssymb,mathtools}
\usepackage{amsthm}
\usepackage{booktabs,array}
\usepackage{enumitem}
\usepackage{algorithm}
\usepackage{algpseudocode}
\usepackage{float}
\usepackage{microtype}
\usepackage[hidelinks]{hyperref}
\newtheoremstyle{plainn}{}{}{\itshape}{}{\bfseries}{.}{.5em}%
  {\thmname{#1}\thmnumber{ #2}\thmnote{ (#3)}}
\theoremstyle{plainn}

\newtheorem{thm}{Theorem}
\newtheorem{prop}[thm]{Proposition}
\newtheorem{lem}[thm]{Lemma}
\newtheorem{cor}[thm]{Corollary}
\newtheorem{claim}[thm]{Claim}
\newtheorem{defn}{Definition}
\theoremstyle{remark}
\newtheorem{remark}{Remark}
\renewcommand{\qedsymbol}{QED}
\newcommand{\row}{\operatorname{row}}
\newcommand{\rank}{\operatorname{rank}}
\newcommand{\Ann}{\operatorname{Ann}}
\newcommand{\Gr}{\operatorname{Gr}}
\newcommand{\Prob}{\operatorname{Prob}}
\newcommand{\poly}{\operatorname{poly}}
\setlength{\emergencystretch}{1.5em}

\begin{document}

\title{Image-Space Reduction and Polynomial-Time Feasibility for Functional Observers
with a Two-Dimensional Compressed Kernel}
\author{}
\maketitle

\begin{abstract}
For a linear system $(A,C)$ and a prescribed observer row block $L_0$, feasibility of a functional
observer of order $q$ is governed by two rank conditions on the pencil $\lambda H-H_A$, where
$H=\mathcal L N$, $H_A=\mathcal L AN$ and $N$ spans the unobservable subspace $K_0=\ker[C;CA]$. The layer index of a design
is $c=k-\rank H$, $k=\dim K_0$; the case $c=1$ was settled in earlier work, where the admissible kernel
directions are the real eigenvectors of the compression $M=\pi_{K_0}A|_{K_0}$.

This paper treats $c=2$. Our main result is a \emph{reduction of the fibre variable}: for a fixed kernel
direction $D$ the pencil depends on the fibre $W$ only through its image
$\Phi(W)\subseteq D^\circ$; the pencil's whole row space lies in $D^\circ$, of dimension $\omega$, but
$\Phi(W)$ itself lies in the image $\varphi(T_D)$, whose dimension is at most $n-k$ (at most twice the
number of output rows). Consequently the search over a Grassmannian of arbitrary dimension collapses to
a search over subspaces of $\varphi(T_D)$. We
give the exact realizability conditions for a prescribed $\Phi$, the resulting reduced feasibility
criterion, a classification of the invariant-plane base for simple spectrum --- with the real and the
conjugate-complex branches treated separately --- a two-valued decision algorithm, and a
$\mathrm{poly}(n,\tau)$ bit-complexity bound under simple spectrum. We also show that the layer index is
\emph{not monotone} at fixed order: there are instances where the $c=2$ layer is feasible while the
$c=1$ layer is not. All structural claims are stated for the split problem; the unrestricted problem is
left open.
\end{abstract}

\begin{IEEEkeywords}
Functional observers, minimum order, compressed pencil, invariant subspaces, Grassmann bundles,
subresultants, bit complexity.
\end{IEEEkeywords}

\section{Introduction}\label{sec:intro}

This section fixes the problem, records what the companion papers settle and what they leave open,
lists the contributions, and closes with the notation used throughout.

\emph{Setting.} The pair $(A,C)$ is given, together with a prescribed row block $L_0$. A design of order
$q$ is a row block $\mathcal L=[L_0;R]$ with $q$ rows; its row space $V$ has dimension $q$. Write
$K_0=\ker[C;CA]$, $k=\dim K_0$, fix a basis $N$ of $K_0$, and put $H=\mathcal L N$, $H_A=\mathcal L AN$. The design is
\emph{feasible} when
\begin{equation}\label{eq:5a}
\row H_A\subseteq\row H,
\end{equation}
\begin{equation}\label{eq:5b}
\rank(\lambda H-H_A)=\omega\quad\text{for all }\lambda\in\mathbb C,
\end{equation}
with $\omega=\rank H$. Condition~\eqref{eq:5b} is equivalent to the vanishing of the gcd of all
$\omega\times\omega$ maximal minors of the pencil (a standard polarity/resultant fact; see
Section~\ref{sec:admissible} and \cite{vzgg13}).

\emph{Prior work and what is new.} The existence question, the order spectrum and the minimum-order
characterisation were settled in \cite{companion1}, \cite{companion2} (see also \cite{fer26a},
\cite{fer26b}): for every order $q$ there is an exact criterion, and the feasible orders are monotone and
saturate. The \emph{complexity} of the general decision problem remains open. In \cite{companion3} the
case $c=1$ was solved structurally: the admissible kernel directions are the real eigenvectors of the
compression $M=\pi_{K_0}A|_{K_0}$, restricted by a prescribed-row orthogonality condition; the design
family is a finite union of relative Grassmann bundles; and one function-field gcd per chart decides the
whole chart.

\emph{Contributions.} This paper treats $c=2$. The contributions are:
\begin{enumerate}[label=\arabic*.,leftmargin=*,itemsep=1pt,topsep=2pt]
\item \emph{Layer bookkeeping (Section~\ref{sec:layers}).} The layer index is $c=k-\omega$ and the
``ambient row excess'' is $t=q-\omega$; the identity $c-t=k-q$ shows that $c$ and $t$ move together when
$q$ is fixed. A previous, ambiguous identification of $\omega$ with $\dim(V\cap K_0)$ is corrected:
$\omega=\rank H=k-\dim(K_0\cap V^\perp)$, while $\dim(V\cap K_0)=\omega$ holds only under the split
condition. We
separate the unrestricted problem \textbf{U} from the $K_0$-split problem \textbf{S} and show by a
counterexample that the split hypothesis cannot be derived from \eqref{eq:5a}.
\item \emph{The two fixed-kernel fibres (Section~\ref{sec:fibres}).} For fixed $D=\ker H$ we give the
exact fibre in both problems, including the exactness condition $K_0\cap(D+AD)=D$ in \textbf{U}, and the
\emph{two independent} prescribed-row conditions in \textbf{S} --- one on the base
($P_{K_0}(V_0)\subseteq D^{\perp_{K_0}}$) and one on the fibre ($W_0\subseteq W$).
\item \emph{Image-space reduction (Section~\ref{sec:reduction}).} The pencil's row space equals the row
space of its $D$-part plus $\Phi(W)\subseteq D^\circ$; since the whole pencil lies in $D^\circ$ (of
dimension $\omega$) while $\Phi(W)$ lies in $\varphi(T_D)$ (of dimension at most $n-k$), the fibre acts
only through a subspace of the low-dimensional image $\varphi(T_D)$. We give the general proof, the exact
realizability inequalities for a prescribed $\Phi$, and the reduced feasibility criterion.
\item \emph{A two-valued algorithm at $c=2$ (Sections~\ref{sec:base}--\ref{sec:decision}).} For simple
spectrum the invariant-plane base is finite, of $O(k^2)$ points, split into a real-real and a
conjugate-complex branch. The image space has at most one projective-line component, and on an
chart every real point is admissible, so the generic-gcd criterion is sharp in both
directions and the algorithm decides exactly.
\item \emph{Bit complexity (Section~\ref{sec:complexity}).} Under simple spectrum the decision runs in
$\poly(n,\tau)$ bit operations, with arithmetic in number fields of degree at most $k^2$.
\item \emph{A non-monotonicity example (Section~\ref{sec:layers:c}).} At one and the same order, the
$c=2$ layer can be feasible while the $c=1$ layer is not. This rules out ``only small layers matter'' as
a complexity strategy.
\end{enumerate}

\emph{Approach and structure.} The approach is uniform: a statement about the fibre $W$ is first pushed
to its image $\Phi(W)$, and the resulting condition is then read on finitely many charts, one per
subspace of $D^\circ$. Section~\ref{sec:layers} fixes the layer parameters and the split hypothesis.
Section~\ref{sec:fibres} computes the fixed-kernel fibres in the two problems.
Section~\ref{sec:reduction} proves the image-space reduction and the reduced criterion.
Section~\ref{sec:base} classifies the base at $c=2$ for simple spectrum and records the general $c$.
Section~\ref{sec:decision} gives the algorithm, its correctness and its bit cost, and closes with a
worked example (Section~\ref{sec:example}).
Section~\ref{sec:conclusion} consolidates the limitations and the evidence.
Appendix~\ref{app:height} holds the height-propagation chain used in the complexity bound.

\emph{Notation.} $E=\mathbb R^n$ throughout, with the Euclidean inner product. $\pi_{K_0}$ is the
orthogonal projection, $K_0^\perp$ its kernel, and $K_0^*$ the dual space of $K_0$. For a subspace
$J\le K_0$ we write
\[
J^\circ:=\{\varphi\in K_0^*:\ \varphi|_J=0\}\le K_0^*
\]
for its annihilator in $K_0^*$, so that $\dim J^\circ=k-\dim J$; the case used
throughout is $J=D=\ker H$, for which $\dim D=c=k-\omega$ and hence $\dim D^\circ=k-c=\omega$. $\Gr(d,X)$ is the
Grassmannian of $d$-planes in $X$. All matrices are real; $p$ is the number of output rows and $\tau$ the
maximum bit length of the entries of $A$ and $C$. For a subspace $X$ we write $\dim X$ for its dimension.

\section{The layer parameters}\label{sec:layers}

This section fixes the three integer parameters attached to a design, isolates the split hypothesis under
which the structural results are stated, and records a counterexample showing that the layer index is not
monotone in the order. Parts A and B are bookkeeping; Part C is the reason the later sections treat $c$
as given rather than assuming the smallest layer.

\subsection{Three quantities that must not be confused}\label{sec:layers:a}

For a design $\mathcal L$ with row space $V$ of dimension $q$, define
\begin{equation}\label{eq:params}
\omega:=\rank H,\qquad c:=k-\omega,\qquad t:=q-\omega .
\end{equation}

\begin{claim}[layer identities]\label{clm:identities}
Let $D=\ker H=K_0\cap V^\perp$. Then
\begin{enumerate}[label=(\arabic*),leftmargin=*,itemsep=1pt,topsep=2pt]
\item {\raggedright $\omega=\rank H=\rank(\mathcal L N)=\dim\pi_{K_0}(V)
=k-\dim(K_0\cap V^\perp)=k-\dim D$;\par}
\item $c=k-\omega=\dim D$;
\item $t=q-\omega=\dim(V\cap K_0^\perp)$, and $q=\omega+t$;
\item in the split case $\dim(V\cap K_0)=\omega$, and then
$V=(V\cap K_0)\oplus(V\cap K_0^\perp)$; in general $\dim(V\cap K_0)$ is a third quantity, not determined
by $\omega$ and $t$ (Claim~\ref{clm:notnormal} below has $\dim(V\cap K_0)=1$ with $c=2$ and $t=3$);
\item $\dim D^\circ=k-\dim D=k-c=\omega$.
\end{enumerate}
Moreover $c-t=k-q$.
\end{claim}

\begin{proof}
For (1): $\ker(\mathcal L N)=\{u\in K_0:\mathcal L u=0\}$, i.e.\
$K_0\cap\ker\mathcal L=K_0\cap V^\perp$ (when $\mathcal L$ has full row rank, $\ker\mathcal L=V^\perp$);
hence $\omega=\rank H=k-\dim(K_0\cap V^\perp)$. On the other hand, the restriction map
$r:V\to K_0^*$, $v\mapsto\langle v,\cdot\rangle|_{K_0}$, has kernel $V\cap K_0^\perp$, so
$\rank r=q-\dim(V\cap K_0^\perp)=\dim\pi_{K_0}(V)$. Its adjoint $r^*:K_0\to V^*$,
$x\mapsto\langle x,\cdot\rangle|_V$, is the transpose of $\pi_{K_0}|_V$ and has kernel $K_0\cap V^\perp$;
hence $\dim\pi_{K_0}(V)=\rank r=\rank r^*=k-\dim(K_0\cap V^\perp)$. All these
quantities coincide. Part (2) is the definition of $c$ together with (1). Part (3) is the first
computation above, and $q=\omega+t$ is the definition of $t$. Part (4) is the definition of the split
condition (Definition~\ref{defn:split}); the verification of the counterexample's numbers is part of the
proof of Claim~\ref{clm:notnormal}. Part (5) is $\dim D^\circ=k-\dim D=k-c=\omega$, and
$c-t=(k-\omega)-(q-\omega)=k-q$.
\end{proof}

\begin{remark}\label{rem:hyp}
In the hyperplane layer $q=k$ one has $c=t=k-\omega$, so $c$ and $t$ coincide there; this is why the
distinction between them was invisible in \cite{companion3}. It does not make $\omega$ coincide with
either: $\omega=k-c$, and $\dim(V\cap K_0)=\omega$ holds only in the split case. For $q\ne k$ even $c$
and $t$ differ, by $c-t=k-q$.
\end{remark}

Thus $c$ and $t$ coincide only in the hyperplane layer $q=k$; outside it $c\ne t$ by $c-t=k-q$. The layer
is indexed by $c=k-\omega$, not by $\dim(V\cap K_0)$, and $\dim D^\circ=\omega$ is a different quantity
from both.

\subsection{The split hypothesis}\label{sec:layers:b}

\begin{defn}\label{defn:split}
A design is \textbf{$K_0$-split} (problem \textbf{S}) when $\dim(V\cap K_0)=\omega$; equivalently
$V=(V\cap K_0)\oplus(V\cap K_0^\perp)$. Designs satisfying \eqref{eq:5a}, \eqref{eq:5b} without this
condition form the \textbf{unrestricted} problem \textbf{U}.
\end{defn}

For $c\ge0$ put $\mathcal I_c(M):=\{D\le K_0:\ \dim D=c,\ M(D)\subseteq D\}$; in the case $c=2$ we write
$\mathcal I_2(M)=\mathcal I_c(M)$.

\begin{claim}[invariant kernel under split]\label{clm:invariant}
Let $\mathcal L$ be split and satisfy \eqref{eq:5a}, and put $D=\ker H=K_0\cap V^\perp$,
$M=\pi_{K_0}A|_{K_0}$. Then $M(D)\subseteq D$.
\end{claim}

\begin{proof}
By \eqref{eq:5a}, $\ker H\subseteq\ker H_A$, i.e.\ $AD\perp V$. Let $\varphi_d(u)=\langle u,d\rangle$ and
$\varphi_{Ad}(u)=\langle u,Ad\rangle$ for $u\in K_0$. Both vanish on $V\cap K_0$: the first because
$D\subseteq V^\perp$, the second by $V\perp AD$. Since the design is split, the annihilator
$(V\cap K_0)^\circ\subseteq K_0^*$ has dimension $k-\dim(V\cap K_0)=k-\omega=c$, and the $\varphi_d$,
$d\in D$, are $c$ independent elements of it; hence $\varphi_{Ad}\in\operatorname{span}\{\varphi_d\}$ and
$\pi_{K_0}(Ad)\in D$.
\end{proof}

\begin{claim}[the split hypothesis is not a normalisation]\label{clm:notnormal}
There is a design satisfying \eqref{eq:5a}, with exact kernel and $c=2$, whose $D=\ker H$ is not
$M$-invariant; the design is then not $K_0$-split.
\end{claim}

\begin{proof}[Construction]
Let $e_1,\dots,e_8$ be the standard basis of $E=\mathbb R^8$ and put
\[
K_0:=\operatorname{span}\{e_5,\dots,e_8\},\qquad f_i:=e_{4+i}\quad(i=1,\dots,4),
\]
so that $f_1,\dots,f_4$ is a basis of $K_0$ and $e_1,\dots,e_4$ is a basis of $K_0^\perp$. All vectors of
$K_0$ below are written in the $f$-basis. Let
\[
D:=\operatorname{span}\{f_1+f_2,\ f_3+f_4\}\le K_0,\qquad \dim D=2 .
\]
Fix the compression $M=\operatorname{diag}(1,1,2,3)$ in the $f$-basis and the coupling $A_{21}$ below,
whose rows are indexed by $e_1,\dots,e_4$ and whose columns by $f_1,\dots,f_4$:
\[
A_{21}:=\begin{pmatrix}
1&-1&0&0\\
0&0&1&0\\
0&0&0&0\\
0&0&0&0
\end{pmatrix},
\]
and define $A$ on the basis $(f_1,\dots,f_4,e_1,\dots,e_4)$ of $E$ by
\begin{equation}\label{eq:notnormal-A}
\begin{gathered}
A=\begin{pmatrix} A_{11} & A_{12}\\ A_{21} & A_{22}\end{pmatrix}
=\begin{pmatrix} M & 0\\ A_{21} & 0\end{pmatrix},\\
Af_i=Mf_i+A_{21}f_i,\qquad Ae_j=0\quad(j=1,\dots,4).
\end{gathered}
\end{equation}
Thus $A_{11}=M=\pi_{K_0}A|_{K_0}$ is the compression, $A_{21}=P_{K_0^\perp}A|_{K_0}$ the coupling to
$K_0^\perp$, and $A_{12}=A_{22}=0$. For $d_1:=f_1+f_2$ and $d_2:=f_3+f_4$ one has
\begin{gather*}
Ad_1=d_1,\\
Ad_2=Md_2+A_{21}d_2=2f_3+3f_4+e_2=2d_2+(f_4+e_2),
\end{gather*}
because $A_{21}d_1=0$ and $A_{21}d_2=e_2$. Hence
\[
U:=D+AD=\operatorname{span}\{d_1,\ d_2,\ f_4+e_2\},\qquad \dim U=3 ,
\]
the three vectors being independent because $f_4+e_2\notin K_0$ while $d_1,d_2\in K_0$. Every
$\alpha d_1+\beta d_2+\gamma(f_4+e_2)\in K_0$ has $\gamma=0$, so
\[
K_0\cap U=K_0\cap(D+AD)=D .
\]
Put $V:=U^\perp$, so $q=\dim V=8-3=5$, and take $\mathcal L$ to be a matrix whose rows form a basis of
$V$; for $L_0=0$ this is a design with row space $V$.
Then $D=\ker H=K_0\cap V^\perp$, $\omega=\rank H=k-\dim D=2$, $c=k-\omega=2$ and $t=q-\omega=3$; the
kernel is exact because $K_0\cap V^\perp=K_0\cap U=D$. Moreover $AD\subseteq U=V^\perp$, so
$\ker H\subseteq\ker H_A$ and \eqref{eq:5a} holds. Finally
\[
V=U^\perp=\operatorname{span}\{e_1,\ e_3,\ e_4,\ e_6-e_5,\ e_8-e_7-e_2\},
\]
whence $V\cap K_0=\operatorname{span}\{e_6-e_5\}$ has dimension $1\ne\omega=2$: the design is not
$K_0$-split. And $Md_2=2f_3+3f_4=2d_2+f_4\notin D$, so $M(D)\not\subseteq D$.
\end{proof}

The split hypothesis is therefore a genuine restriction on the layer rather than a normalisation;
consequently every structural statement below is labelled \textbf{S}, and the only statement for the
unrestricted problem is Theorem~\ref{thm:ufibre}.

\subsection{The layer index is not monotone}\label{sec:layers:c}

At a fixed order $q$ the identity $c-t=k-q$ shows that $c$ and $t$ move together, so only one order
allows both layers to be compared. Take $k=4$: then $c=1$ forces $q=3+t$ and $c=2$ forces $q=2+t$, and
$q=3$ is the common value ($c=1,t=0$; $c=2,t=1$). For the instance $A,C$ of \cite{companion3}, \S VI,
with the compression replaced by $\operatorname{diag}(1,1,2,3)$:

\begin{thm}[non-monotonicity]\label{thm:nonmono}
At order $q=3$ the $c=2$ layer is feasible and the $c=1$ layer is not.
\end{thm}

\begin{proof}
\emph{$c=2$ side.} With $D=\operatorname{span}\{e_1+e_2,e_3\}$ and
$W=\operatorname{span}\{(0,0,1,1,0,0,0,0)\}$, the design $V=D^{\perp_{K_0}}\oplus W$ has $\dim V=3$,
$\omega=2$, satisfies \eqref{eq:5a} and exactness ($\dim(K_0\cap V^\perp)=2$), and the six nonzero maximal
$2\times2$ minors of the pencil are
\[
-(\lambda-3)(\lambda-1),\ (\lambda-3)(\lambda-1),\ \lambda-1,\ 1-\lambda,\ \lambda-3,\ 3-\lambda,
\]
whose gcd is $1$ (indeed $\tfrac12\bigl((\lambda-1)-(\lambda-3)\bigr)=1$). So \eqref{eq:5b} holds.

\emph{$c=1$ side.} At $q=3$ and $c=1$ one has $t=q-\omega=0$, so $\dim(V\cap K_0^\perp)=0$ and the design
is automatically split; Claim~\ref{clm:invariant} therefore applies. The invariant lines are the lines of
the eigenspace $E_1=\operatorname{span}\{e_1,e_2\}$ together with $\operatorname{span}\{e_3\}$ and
$\operatorname{span}\{e_4\}$. For the family $D=\operatorname{span}\{e_1+se_2\}$ one computes
\[
H=\begin{pmatrix}-s&1&0&0\\0&0&1&0\\0&0&0&1\end{pmatrix},\qquad \rank H=3,
\]
and the gcd of the nonzero maximal minors equals $(\lambda-1)(\lambda-2)(\lambda-3)$,
\textbf{independently of $s$}; the point at infinity $s=\infty$ gives the same gcd, and the two isolated
directions give $(\lambda-1)^2(\lambda-3)$ and $(\lambda-1)^2(\lambda-2)$. Hence no $c=1$ design of order
$3$ is feasible.
\end{proof}

\begin{remark}\label{rem:direction}
The direction of the failure is ``the larger layer is easier'', so no monotonicity argument can replace
the analysis of all layers.
\end{remark}

Consequently the analysis of one layer cannot replace the analysis of another at the same order, and
Sections~\ref{sec:fibres}--\ref{sec:decision} treat $c$ as given.

\section{The fixed-kernel fibre}\label{sec:fibres}

This section computes, for a fixed kernel direction $D$, the exact set of admissible row spaces in the two
problems. Part A treats the unrestricted problem \textbf{U}; Part B the split problem \textbf{S}. In both
cases the answer is a Grassmannian fibre with an explicit dimension and an explicit non-emptiness
condition.

Throughout this section $D\le K_0$ with $\dim D=c$ is fixed, $U_D:=D+AD$, and $V_0=\row L_0$. We write
$P_{K_0}$ and $P_{K_0^\perp}$ for the two orthogonal projections, and $M=P_{K_0}A|_{K_0}$,
$A_{21}=P_{K_0^\perp}A|_{K_0}$, so that $Ax=Mx+A_{21}x$ for $x\in K_0$ (the two summands lie in
complementary subspaces).

\subsection{The unrestricted problem \texorpdfstring{\textbf{U}}{U}}\label{sec:fibres:u}

\begin{thm}[unrestricted fibre]\label{thm:ufibre}
Let $D\le K_0$ with $\dim D=c$.
\begin{enumerate}[label=\arabic*.,leftmargin=*,itemsep=1pt,topsep=2pt]
\item A $q$-dimensional $V$ satisfies $V_0\subseteq V$, $V\subseteq U_D^\perp$ and $K_0\cap V^\perp=D$ if
and only if
\[
V_0\subseteq V\subseteq U_D^\perp,\qquad \dim V=q,\qquad K_0\cap(D+AD)=D .
\]
\item The condition $K_0\cap(D+AD)=D$ is \textbf{equivalent} to exactness of the kernel, i.e. to
$K_0\cap V^\perp=D$ for every $V$ with $V_0\subseteq V\subseteq U_D^\perp$.
\end{enumerate}
\end{thm}

\begin{proof}
(1) If $K_0\cap V^\perp=D$ and \eqref{eq:5a} holds then $U_D\subseteq V^\perp$ and hence
$K_0\cap(D+AD)\subseteq K_0\cap V^\perp=D$, while $D\subseteq D+AD$ gives the reverse inclusion.
Conversely if $K_0\cap(D+AD)=D$, then any $V$ with $V_0\subseteq V\subseteq U_D^\perp$ satisfies
$D\subseteq K_0\cap V^\perp\subseteq K_0\cap(D+AD)=D$. (2) is the same two lines.
\end{proof}

\begin{remark}\label{rem:ufibre}
For $c=2$ the fibre of \textbf{U} is thus
$\{V\in\Gr(q,U_D^\perp):V_0\subseteq V,\ K_0\cap V^\perp=D\}$, which is an open subset of a Grassmannian
once the exactness condition holds. We do not eliminate the fibre variable in \textbf{U}.
\end{remark}

The unrestricted fibre is thus an open subset of a Grassmannian once the exactness condition
$K_0\cap(D+AD)=D$ holds, and that condition is a property of $D$ alone.

\subsection{The split problem \texorpdfstring{\textbf{S}}{S}}\label{sec:fibres:s}

\begin{thm}[split fibre]\label{thm:splitfibre}
Let $D\in\mathcal I_c(M)$ with $\dim D=c$, and put
\[
T_D:=K_0^\perp\cap(A_{21}D)^\perp,\qquad W_0:=P_{K_0^\perp}(V_0),
\]
\[
w_0:=\dim W_0,\qquad r(D):=\rank(A_{21}|_D).
\]
Assume the two prescribed-row conditions
\begin{equation}\label{eq:splithyp}
\textbf{(H1)}\quad P_{K_0}(V_0)\subseteq D^{\perp_{K_0}},
\qquad
\textbf{(H2)}\quad W_0\subseteq T_D .
\end{equation}
Then the row spaces of the $K_0$-split designs with kernel $D$, prescribed rows $V_0$, and satisfying
\eqref{eq:5a} are exactly
\[
V=D^{\perp_{K_0}}\oplus W,\qquad W_0\subseteq W\subseteq T_D,\quad \dim W=t,
\]
and the fibre is
\begin{equation}\label{eq:splitfibre}
\begin{gathered}
\mathcal F_D^{\rm split}\cong\Gr\bigl(t-w_0,\ T_D/W_0\bigr),\\
\dim\mathcal F_D^{\rm split}=(t-w_0)\bigl((n-k)-r(D)-t\bigr),
\end{gathered}
\end{equation}
which is non-empty if and only if
\[
w_0\le t\le (n-k)-r(D).
\]
\end{thm}

\begin{proof}
(i) For a split design, $\dim(V\cap K_0)=\omega=k-c$ by Claim~\ref{clm:identities}(4). Every
$x\in V\cap K_0$ is orthogonal to $D\subseteq V^\perp$, and $\dim D^{\perp_{K_0}}=k-\dim D=k-c$; hence
$V\cap K_0=D^{\perp_{K_0}}$, and $V=(V\cap K_0)\oplus W$ with $W:=V\cap K_0^\perp$. Then
$K_0\cap V^\perp=D$: the inclusion $D\subseteq K_0\cap V^\perp$ holds because $D\perp D^{\perp_{K_0}}$ and
$D\perp W$, and the reverse inclusion follows from
$K_0\cap V^\perp\subseteq K_0\cap(D^{\perp_{K_0}})^\perp=D$.
(ii) Recall that \eqref{eq:5a} is $\ker H\subseteq\ker H_A$, i.e.\ $AD\perp V$: indeed
$\ker H_A=\{u\in K_0:Au\in V^\perp\}$. For $d\in D$, $\langle d,V\rangle=0$; writing $Ad=Md+A_{21}d$
with $Md\in D\subseteq V^\perp$ and $A_{21}d\in K_0^\perp$, the relation $\langle Ad,V\rangle=0$ is
equivalent to $W\perp A_{21}D$. Hence \eqref{eq:5a} holds if and only if $W\subseteq T_D$, i.e.\ if and
only if the fibre parameter lies in $T_D$.
(iii) $V_0\subseteq V$ splits into $P_{K_0}(V_0)\subseteq D^{\perp_{K_0}}$ and $W_0\subseteq W$, which are
(H1) and the lower inclusion of the fibre.
(iv) Since $A_{21}D\subseteq K_0^\perp$ and $\dim A_{21}D=r(D)$, its orthogonal complement inside
$K_0^\perp$ has dimension $(n-k)-r(D)$; thus $\dim T_D=(n-k)-r(D)$, and $W/W_0$ is a free
$(t-w_0)$-plane in $T_D/W_0$, a Grassmannian of dimension
$(t-w_0)(\dim T_D-t)=(t-w_0)((n-k)-r(D)-t)$.
\end{proof}

\begin{remark}[the prescribed rows are two conditions]\label{rem:presrows}
The condition on the base,
\begin{equation}\label{eq:bpres}
\begin{aligned}
\mathcal B_{\rm pres}&:=\bigl\{D\in\mathcal I_c(M):\ P_{K_0}(V_0)\subseteq D^{\perp_{K_0}}\bigr\}\\
&=\bigl\{D\in\mathcal I_c(M):\ D\perp P_{K_0}(V_0)\bigr\},
\end{aligned}
\end{equation}
is not implied by $W_0\subseteq W$; omitting it produces ``designs'' that cannot be completed to a legal
observer. It also removes entire components of the base: if $P_{K_0}(V_0)$ contains a vector with
non-zero inner product with every element of a given invariant $D$, that $D$ is excluded regardless of
the fibre.
\end{remark}

The split fibre is thus a Grassmannian of computable dimension over the field of definition of $D$, and
its two prescribed-row conditions are independent; this is the description that
Section~\ref{sec:reduction} compresses further.

\section{Image-space reduction}\label{sec:reduction}

This section proves that, for a split design with $M(D)\subseteq D$, the pencil sees the fibre $W$ only
through its image $\Phi(W)\subseteq D^\circ$. Part A proves the reduction, Part B gives the exact
realisability conditions for a prescribed image, and Part C states the resulting criterion.

Fix $D\in\mathcal B_{\rm pres}$ and a split design $V=D^{\perp_{K_0}}\oplus W$ with
$W_0\subseteq W\subseteq T_D$.

\subsection{The reduction}\label{sec:reduction:a}

Write $H_K,(H_A)_K$ for the rows of $H,H_A$ carried by $D^{\perp_{K_0}}$, and put
$P_D(\lambda):=\lambda H_K-(H_A)_K$, $R_\lambda:=\row P_D(\lambda)$.

Let
\begin{equation}\label{eq:iotaN}
\iota_N:K_0^*\longrightarrow\mathbb R^{1\times k},\qquad \ell\longmapsto \ell N
\end{equation}
be the coordinate isomorphism induced by the basis $N$: the row vector $\ell N$ has entries
$\ell(Ne_j)$, $j=1,\dots,k$, and $\iota_N$ is a linear isomorphism because $N$ has rank $k$. We identify
$K_0^*$ with $\mathbb R^{1\times k}$ through $\iota_N$; from now on every row-space identity is understood
after this identification, and a subspace of $K_0^*$ and its image under $\iota_N$ are written with the
same symbol.

For $v\in E$ let $h_v,h^A_v\in K_0^*$ be the restrictions to $K_0$ of $u\mapsto\langle v,u\rangle$ and
$u\mapsto\langle v,Au\rangle$, and put $\Psi_\lambda(v):=\lambda h_v-h^A_v$. The $j$-th row of
$\lambda H-H_A$ is $\iota_N(\Psi_\lambda(v_j))$, where $v_j$ is the $j$-th row of $\mathcal L$; hence
$\row(\lambda H-H_A)=\iota_N(\Psi_\lambda(V))$.

For $w\in K_0^\perp$ define $\varphi_w\in K_0^*$ by $\varphi_w(u)=\langle w,A_{21}u\rangle$, and let
\[
\varphi:K_0^\perp\longrightarrow K_0^*,\qquad \varphi(w):=\varphi_w
\]
be the induced linear map. For $W\le K_0^\perp$ put
\begin{equation}\label{eq:Phi}
\Phi(W):=\varphi(W)=\{\varphi_w:\ w\in W\}\le D^\circ .
\end{equation}
The inclusion is an inclusion of subspaces, and $\Phi(W)$ is the image of the linear map $\varphi$, not
merely a set: $W\subseteq T_D$ means $W\perp A_{21}D$, hence $\varphi_w(d)=\langle w,A_{21}d\rangle=0$ for
all $d\in D$, so $\varphi(W)\subseteq D^\circ$.

\begin{thm}[image-space reduction]\label{thm:reduction}
For every $\lambda$,
\begin{equation}\label{eq:reduction}
\row\bigl(\lambda H-H_A\bigr)=R_\lambda+\Phi(W),
\end{equation}
the identity holding in $\mathbb R^{1\times k}$ under $\iota_N$; equivalently,
$\Psi_\lambda(V)=\Psi_\lambda(D^{\perp_{K_0}})+\Phi(W)$ in $K_0^*$. Consequently the pencil condition
depends on $W$ only through the subspace $\Phi(W)$, and
$\dim\Phi(W)\le\dim\varphi(T_D)\le\dim T_D=(n-k)-r(D)$. Moreover:
\begin{enumerate}[label=\arabic*.,start=3,leftmargin=*,itemsep=1pt,topsep=2pt]
\item $\row(\lambda H-H_A)\subseteq D^\circ$ for every $\lambda$, and $\dim D^\circ=k-c=\omega$;
\item hence $\rank(\lambda H-H_A)\le\omega$ automatically, and \eqref{eq:5b} is equivalent to
\[
\row\bigl(\lambda H-H_A\bigr)=D^\circ\quad\text{for all }\lambda,
\]
\[
\text{i.e.}\qquad R_\lambda+\Phi(W)=D^\circ\quad\forall\lambda.
\]
\end{enumerate}
\end{thm}

\begin{proof}
\begin{enumerate}[label=\arabic*.,leftmargin=*,itemsep=1pt,topsep=2pt]
\item For $w\in W\le K_0^\perp$ and $u\in K_0$: $h_w(u)=\langle w,u\rangle=0$, and
$h^A_w(u)=\langle w,Au\rangle=\langle w,Mu\rangle+\langle w,A_{21}u\rangle=0+\varphi_w(u)$; hence
$\Psi_\lambda(w)=-\varphi_w$, independently of $\lambda$.
\item For any linear $\Psi$ and any direct sum $V=P\oplus W$, $\Psi(V)=\Psi(P)+\Psi(W)$; here
$P=D^{\perp_{K_0}}$.
\end{enumerate}
Combining, $\Psi_\lambda(V)=\Psi_\lambda(P)+\Phi(W)$, and the sign is immaterial because a subspace
equals its negative; applying the isomorphism $\iota_N$ gives \eqref{eq:reduction}.
\begin{enumerate}[label=\arabic*.,start=3,leftmargin=*,itemsep=1pt,topsep=2pt]
\item For $v\in P=D^{\perp_{K_0}}$ and $d\in D$: $\langle v,d\rangle=0$, and
$\langle v,Ad\rangle=\langle v,Md\rangle+\langle v,A_{21}d\rangle=0+0$ because $Md\in D$ and
$A_{21}d\in K_0^\perp$; hence every $\Psi_\lambda(v)$ annihilates $D$, i.e.\ lies in $D^\circ$, and
$\Phi(W)\subseteq D^\circ$ by \eqref{eq:Phi}. Since $\dim D^\circ=k-c=\omega$, the rank bound follows.
\item A subspace of the $\omega$-dimensional $D^\circ$ has dimension $\omega$ if and only if it is
$D^\circ$.
\end{enumerate}
\end{proof}

\begin{remark}[what the reduction says]\label{rem:whatreduction}
The spectral condition does not see $W$, only $\Phi(W)$: if $\Phi(W_1)=\Phi(W_2)$ then the two designs
have identical pencils up to row operations. The design space is therefore quotiented by the fibres of
$\Phi$, and the relevant parameter lives in the image $\varphi(T_D)$, a space of dimension at most $n-k$.
For $c=2$ this is decisive (Section~\ref{sec:decision}).
\end{remark}

So the spectral condition is a condition on the image $\Phi(W)$ and not on $W$, and the fibre variable
has been replaced by a subspace of $\varphi(T_D)$, a space whose dimension is at most the codimension of
$K_0$.

\subsection{Exact realizability}\label{sec:reduction:b}

\begin{lem}[subspaces with a prescribed image]\label{lem:prescribed}
Let $f:X\to Y$ be linear, let $Z\le\operatorname{im}f$ and let $d\ge0$. There is a subspace $S\le X$
with $f(S)=Z$ and $\dim S=d$ if and only if
\begin{equation}\label{eq:prescribed}
\dim Z\le d\le\dim\ker f+\dim Z .
\end{equation}
\end{lem}

\begin{proof}
If $f(S)=Z$ then $\dim Z=\dim S-\dim(S\cap\ker f)\le d$, and $S\le f^{-1}(Z)$ gives
$d\le\dim f^{-1}(Z)=\dim\ker f+\dim Z$. Conversely assume \eqref{eq:prescribed}. Choose a complement
$S_0$ of $\ker f$ in $f^{-1}(Z)$; then $f|_{S_0}:S_0\to Z$ is an isomorphism, so $\dim S_0=\dim Z$.
Since $\dim f^{-1}(Z)=\dim\ker f+\dim Z\ge d$, extend $S_0$ inside $f^{-1}(Z)$ to a subspace $S$ of
dimension $d$. Then $f(S)=f(S_0)=Z$.
\end{proof}

\begin{thm}[realizability of a prescribed image]\label{thm:realize}
Assume (H2) of \eqref{eq:splithyp}, so that $T_D/W_0$ is defined. Let $\Phi_0:=\varphi(W_0)$ and let
$\bar\varphi:T_D/W_0\to D^\circ/\Phi_0$ be the map induced by
$\varphi|_{T_D}$. A subspace $\Phi\le D^\circ$ is realised by
some $W$ with $W_0\subseteq W\subseteq T_D$ and $\dim W=t$ if and only if
\begin{equation}\label{eq:realize}
\begin{aligned}
&\Phi_0\subseteq\Phi\subseteq\varphi(T_D),\\
&\dim\Phi-\dim\Phi_0\ \le\ t-w_0\\
&\qquad\le\ \dim\ker\bar\varphi+\dim\Phi-\dim\Phi_0 .
\end{aligned}
\end{equation}
\end{thm}

\begin{proof}
Subspaces $W$ with $W_0\subseteq W\subseteq T_D$ correspond bijectively to subspaces
$S=W/W_0\le T_D/W_0$, with $\dim S=\dim W-w_0$, and $\varphi(W)=\Phi$ if and only if
$\bar\varphi(S)=\Phi/\Phi_0$ (which presupposes $\Phi_0\subseteq\Phi$). Apply
Lemma~\ref{lem:prescribed} to $f=\bar\varphi$, $Z=\Phi/\Phi_0$ and $d=t-w_0$: an $S$ with
$\bar\varphi(S)=Z$ and $\dim S=d$ exists if and only if $Z\le\operatorname{im}\bar\varphi$, i.e.\
$\Phi\subseteq\varphi(T_D)$, and
$\dim Z\le d\le\dim\ker\bar\varphi+\dim Z$, which is \eqref{eq:realize}.
\end{proof}

Every subspace of $D^\circ$ between $\Phi_0$ and $\varphi(T_D)$ whose dimension lies in the stated
interval is therefore realised, and no other subspace is; the two inequalities are what the enumeration
of Part C quantifies over.

\subsection{The reduced criterion}\label{sec:reduction:c}

\begin{thm}\label{thm:reduced}
Recall $P_D(\lambda)=\lambda H_K-(H_A)_K$ and $r_D(\lambda)=\dim R_\lambda$. Then
\begin{align}
&\exists W\ (W_0\subseteq W\subseteq T_D,\ \dim W=t)\ \text{satisfying \eqref{eq:5b}}\notag\\
&\iff\ \exists\Phi\ \text{satisfying Theorem~\ref{thm:realize} such that}\notag\\
&\qquad\dim\bigl(\row P_D(\lambda)+\Phi\bigr)=\omega\quad\forall\lambda,
\label{eq:reduced}
\end{align}
and the latter is equivalent to
\begin{equation}\label{eq:reduced2}
\dim\Phi-\dim\bigl(\Phi\cap\row P_D(\lambda)\bigr)=\omega-r_D(\lambda)\quad\forall\lambda .
\end{equation}
\end{thm}

\begin{proof}
Immediate from Theorem~\ref{thm:reduction} and $\dim(X+Y)=\dim X+\dim Y-\dim(X\cap Y)$.
\end{proof}

\begin{remark}\label{rem:imagesize}
Theorems~\ref{thm:reduction}--\ref{thm:reduced} reduce the fibre problem to a quantifier over subspaces of
$D^\circ$. Two spaces matter and they must not be confused: $D^\circ$ has dimension $\omega$, which grows
with $k$, while the image $\varphi(T_D)$ has dimension at most
$\dim T_D=(n-k)-r(D)\le n-k\le 2p$, where $p$ is the number of output rows. So the effective parameter
ranges over
\[
\bigcup_{j\le\min(t,\,2p)}\Gr\bigl(j,\ \varphi(T_D)\bigr),
\]
a family of dimension at most $\lfloor\min(t,2p)^2/4\rfloor\le p^2$, \textbf{independent of $c$}. In
particular the image parameter is a point, a projective line, or a plane only when
$\dim\varphi(T_D)\le2$; for larger output counts the parameter is a higher-dimensional Grassmannian. A
necessary condition for feasibility follows at once:
$\max_\lambda(\omega-r_D(\lambda))\le\dim\Phi\le n-k$, i.e. the rank deficiency of $P_D(\lambda)$ can
never exceed the codimension of $K_0$.
\end{remark}

\section{The base at \texorpdfstring{$c=2$}{c=2}}\label{sec:base}

This section describes the base of invariant two-planes at $c=2$ for simple spectrum, and then records
that the reduction of Section~\ref{sec:reduction} is verbatim for arbitrary $c$. Part A is the
classification, Part B the general statement.

Write $\mathcal I_2(M)$ for $\mathcal I_c(M)$ with $c=2$ (Section~\ref{sec:layers}).

\subsection{Simple spectrum: the invariant-plane base}\label{sec:base:a}

\begin{thm}[base, simple spectrum]\label{thm:base}
Suppose $M$ has simple spectrum, i.e. all eigenvalues of $M$ over $\mathbb C$ are distinct, equivalently the characteristic polynomial of $M$ is squarefree. Then $\mathcal I_2(M)$ is finite, and every element is one of:
\begin{enumerate}[label=\arabic*.,leftmargin=*,itemsep=1pt,topsep=2pt]
\item \textbf{(branch R)} $\operatorname{span}\{v_\alpha,v_\beta\}$ for two real eigenvalues
$\alpha\ne\beta$ with eigenvectors $v_\alpha,v_\beta$;
\item \textbf{(branch C)} the real two-plane attached to a complex-conjugate pair: for a real
eigenvector decomposition $z=u+iw$ of an eigenvalue $\mu=a+ib$ with $b\ne0$, the plane
$\operatorname{span}\{u,w\}$, equivalently $\ker f(M)$ where $f$ is the irreducible quadratic factor.
\end{enumerate}
The total number of planes is $\binom{k_{\mathbb R}}{2}+\#\{\text{conjugate pairs}\}\le k^2$, where
$k_{\mathbb R}$ is the number of real eigenvalues.
\end{thm}

\begin{proof}
For a simple spectrum $M$ is diagonalisable over $\mathbb C$ and each complex eigenspace is a line. An
invariant real two-plane either contains two real eigenvectors (branch R), or is spanned by the real and
imaginary parts of a complex eigenvector (branch C), or is $\mathbb R^2$ for a two-dimensional $M$ with a
non-real pair, which is the case $k=2$ of branch C. Nothing else is invariant, because the restriction of
$M$ to $D$ has a characteristic polynomial of degree two dividing the minimal polynomial. The count is
immediate.
\end{proof}

\begin{remark}[what is \emph{not} claimed]\label{rem:nonsimple}
For non-simple spectra $\mathcal I_2(M)$ can be positive-dimensional (the families $\Gr(2,K_\lambda)$
inside a repeated eigenspace, and products $\mathbb P(K_\lambda)\times\mathbb P(K_\mu)$), and the planes
attached to a repeated irreducible quadratic form a $\mathbb{CP}^{g-1}$-family of real dimension
$2(g-1)$, $g$ the geometric multiplicity. We make \textbf{no} claim of a rational parametrisation of the
general invariant-plane variety; the simple-spectrum statement above is what the algorithm uses.
\end{remark}

\begin{table*}[!t]
\centering
\caption{The invariant-plane classification of Theorem~\ref{thm:base} at $c=2$ under simple spectrum:
the two branches, the field over which each plane is defined (Claim~\ref{clm:arith}), and the number of
planes. Here $k_{\mathbb R}$ is the number of real eigenvalues of $M$ and $k=\dim K_0$.}
\label{tab:planes}
\small
\begin{tabular}{@{}l>{\raggedright\arraybackslash}p{0.44\textwidth}>{\raggedright\arraybackslash}p{0.24\textwidth}l@{}}
\toprule
Branch & Plane $D$ & Defined over & Number \\
\midrule
R & $\operatorname{span}\{v_\alpha,v_\beta\}$, two real eigenvalues $\alpha\ne\beta$
& $\mathbb Q(\alpha,\beta)$, degree $\le k^2$ & $\binom{k_{\mathbb R}}{2}$ \\
C & $\operatorname{span}\{u,w\}$ from $z=u+iw$ of $\mu=a+ib$, $b\ne0$; equivalently
$\ker f(M)$, $f$ irreducible quadratic & $\mathbb Q$ & $\#\{\text{conjugate pairs}\}$ \\
\midrule
Total & & & $\le k^2$ \\
\bottomrule
\end{tabular}
\end{table*}

Under simple spectrum the base is therefore finite and explicit, which is what makes the enumeration of
Section~\ref{sec:algorithm} terminate in polynomial time; Remark~\ref{rem:nonsimple} records what is not
claimed without simplicity.

\subsection{The general \texorpdfstring{$c$}{c}}\label{sec:base:b}

\begin{prop}[verbatim generalisation]\label{prop:general}
Theorems~\ref{thm:ufibre}--\ref{thm:reduced} hold with $2$ replaced by $c$: the base is
$\mathcal B_{\rm pres}\subseteq\mathcal I_c(M)$, the fibre is $\Gr(t-w_0,T_D/W_0)$ with
$T_D=K_0^\perp\cap(A_{21}D)^\perp$, the image space satisfies $\Phi(W)\subseteq D^\circ$ with
$\dim D^\circ=k-c=\omega$, and the reduced criterion of Theorem~\ref{thm:reduced} is unchanged. The
effective image parameter ranges over $\bigcup_{j\le\min(t,2p)}\Gr(j,\varphi(T_D))$, of dimension at most
$\lfloor\min(t,2p)^2/4\rfloor\le p^2$.
\end{prop}

\begin{proof}
The proofs of Theorems~\ref{thm:reduction}--\ref{thm:reduced} use only $\dim D=c$ and $M(D)\subseteq D$;
the numerical value $2$ enters only in the list of possible image dimensions.
\end{proof}

\begin{remark}\label{rem:annihilator}
Note that $\dim D^\circ=\omega$, not $c$: since $c=k-\omega$, the annihilator of the $c$-dimensional
subspace $D$ in the $k$-dimensional space $K_0^*$ has dimension $k-c=\omega$. The annihilator of
$V\cap K_0$ is a different space: in the split case it has dimension $k-\dim(V\cap K_0)=c$, and it is
$D^\circ$, not $(V\cap K_0)^\circ$, that contains the pencil. For $c=2$ and $k=4$ the two dimensions
coincide, which is why the small cases are misleading; in general $D^\circ$ is large and only the image
$\varphi(T_D)$ is small.
\end{remark}

\begin{remark}[what remains open]\label{rem:open}
For general $c$ three obstructions persist: the enumeration of $\mathcal I_c(M)$ (exponential when $c$
varies), the choice of $\Phi$ (a Grassmannian of dimension up to $p^2$), and the absence of a rational
parametrisation of the base in the non-simple case. Two questions seem most promising: (a) whether the
quantifier over $\Phi$ can be replaced by a rank or separation condition on the induced maps of
$T_D/W_0$; (b) whether the layer index can be traded for $\dim U$ with $U=\pi_1(W)$, since
$c=k-\dim U$. We do not know the answer to either.
\end{remark}

The reduction of Section~\ref{sec:reduction} therefore does not use $c=2$; what $c=2$ buys is a finite
base (Theorem~\ref{thm:base}), while the smallness of the image parameter in Section~\ref{sec:decision}
comes from the bound on $n-k$, not from the smallness of $c$ (Remark~\ref{rem:boundaries}).

\section{Decision and its cost}\label{sec:decision}

This section assembles the decision procedure at $c=2$: Part A shows that admissibility is automatic on
an chart, Part B states the algorithm, Part C proves it correct, and Part D bounds its bit
cost.

\subsection{Admissibility is automatic on every chart}\label{sec:admissible}

Recall from \cite{companion3} that a point of a chart is \emph{admissible} when some maximal minor of the
cleared pencil keeps its generic $\lambda$-degree at that point; the specialization lemma then shows that
the generic gcd degree is a lower bound on every admissible point, and that the generic gcd being a
constant forces the chart to contain a real feasible point whenever the real points are Zariski dense.

\begin{claim}\label{clm:admissible}
On every chart of a split design, $\rank H=\omega$ identically. In particular $H$ is independent of all
chart parameters, so the $\lambda$-leading coefficient of a suitable maximal minor is a non-zero constant
and every real point of every chart is admissible.
\end{claim}

\begin{proof}
The $W$-rows of $H$ vanish (Theorem~\ref{thm:reduction}, step 1) and the $D$-rows have rank
$\dim D^{\perp_{K_0}}=k-c=\omega$; hence $\rank H=\omega$ for every choice of $W$. On any chart of
$\Gr(j,V_\varphi)$ the rows of $H$ are parameter-free, and the $\lambda$-leading coefficient of the
maximal minor built on those $\omega$ rows is a non-zero constant.
\end{proof}

Admissibility being automatic, the generic-gcd criterion of \cite{companion3} becomes an exact test rather
than a generic one, and the algorithm can be stated in Part B without a certificate per point.

\subsection{The algorithm}\label{sec:algorithm}

\begin{defn}[image charts]\label{defn:charts}
Let $\Phi_0=\varphi(W_0)$ and $V_\varphi:=\varphi(T_D)$, of dimension $\le\min(n-k,2p)$. The realisable
image spaces are the subspaces $\Phi\le V_\varphi$ with $\Phi_0\subseteq\Phi$, $\dim\Phi\le t$, satisfying
the inequalities of Theorem~\ref{thm:realize}; they are charted by the standard affine charts of
$\bigcup_{j}\Gr(j,V_\varphi)$, each of which carries $j(\dim V_\varphi-j)\le p^2$ parameters. Two
degenerate cases deserve names: $\Phi=\{0\}$ (possible only if $\Phi_0=0$ and $t\ge0$ with the lower
inequality satisfied) and $\Phi=D^\circ$ itself (possible only if $V_\varphi=D^\circ$, i.e. only in the
boundary case $\dim T_D=\omega$). When $\dim\Phi_0\ge1$ the one-dimensional families disappear; the line
case occurs only when $\Phi_0=0$ and $\dim V_\varphi\ge2$.
\end{defn}

The enumeration is finite and explicit, so what remains is to say on which charts the answer is decided
and at what cost.

\begin{algorithm}[H]
\caption{Decision at $c=2$, split, simple spectrum.}
\label{alg:decide}
\begin{algorithmic}[1]
\State compute $K_0$, $N$, $M$, $A_{21}$, $\omega$, $c$, $t$;
\State enumerate $\mathcal I_2(M)$ by Theorem~\ref{thm:base}, keeping only
$D\in\mathcal B_{\rm pres}$;
\State for each $D$ compute $r(D)$, $T_D$, $W_0$, $\Phi_0$, and check $W_0\subseteq T_D$;
\State enumerate the realisable image spaces $\Phi$ by Theorem~\ref{thm:realize} and
Remark~\ref{rem:imagesize};
\State for each image chart, clear denominators and compute the gcd of the
$\omega\times\omega$ maximal minors over $\mathbb Q(s)[\lambda]$ (or over the field of definition of
$D$); if it is constant, return \emph{feasible};
\State if every chart gives a non-constant gcd, return \emph{infeasible}.
\end{algorithmic}
\end{algorithm}

\subsection{Correctness}\label{sec:correct}

\begin{claim}[two-valued decision]\label{clm:twovalued}
Under the hypotheses of Theorem~\ref{thm:correct} below, each real-admissible image chart satisfies both
hypotheses of the generic-gcd criterion, hence
\[
\begin{aligned}
&\text{the generic gcd on the chart is }1\\
&\quad\iff\text{the chart contains a real feasible design}.
\end{aligned}
\]
Therefore the algorithm outputs exactly \textbf{feasible} or \textbf{infeasible}.
\end{claim}

\begin{thm}[correctness]\label{thm:correct}
Assume $c=2$, simple spectrum of $M$, and the split condition. Then Algorithm~\ref{alg:decide} terminates
and its answer is correct.
\end{thm}

\begin{proof}
Finite enumeration by Theorems~\ref{thm:base} and~\ref{thm:realize}; correctness on each chart by
Claim~\ref{clm:twovalued} and Theorem~\ref{thm:reduced}; the chart cover is complete because every
realisable $\Phi$ is one of the charts of Definition~\ref{defn:charts}.
\end{proof}

Correctness being established, the only question left is the cost of the enumeration and of the gcd
computations, which is Part D.

\subsection{Bit complexity}\label{sec:complexity}

Let $\tau$ be the maximum bit length of the entries of $A$ and $C$.

\begin{claim}[arithmetic of the enumeration]\label{clm:arith}
Assume simple spectrum. Then the number of planes is $O(k^2)$, each plane is defined over a number field
$L_D$ of degree at most $k^2$, and the field data and the bases of $D$, $D^{\perp_{K_0}}$, $T_D$,
together with $\varphi$, $\Phi_0$ and the entries of the cleared pencil, have heights
$\poly(n,\tau)$.
\end{claim}

\begin{proof}[Sketch]
$M$ is rational, so its minimal polynomial has coefficients of polynomial bit length; its irreducible
factors over $\mathbb Q$ then have polynomial height, and the roots have height bounded by the Mahler
measure of these factors. Branch R uses $L_D=\mathbb Q(\alpha,\beta)$ with
$[\mathbb Q(\alpha,\beta):\mathbb Q]\le d_id_j\le k^2$; eigenvectors are obtained by Cramer's rule over
$L_D$, with denominators $f_i'(\alpha)\ne0$ by simplicity. Branch C uses $\ker f(M)$, which is
$\mathbb Q$-defined, so the final objects live over $\mathbb Q$ even though the intermediate eigenvector
data are computed in a complex embedding of $\mathbb Q[t]/(f)$. Orthogonal complements and kernels are
again Cramer, and the pencil entries are obtained by polynomially many field operations.
\end{proof}

The height propagation is checked step by step in Appendix~\ref{app:height}.

\begin{claim}[degrees on an image chart]\label{clm:degrees}
On a chart of $\Gr(j,V_\varphi)$ at most $j$ rows of the pencil depend on the chart parameters, and each
dependence is affine. Hence every maximal minor satisfies
\[
\deg_{\text{param}}\le j\le\min(t,\,2p),\qquad \deg_\lambda\le\omega\le k .
\]
In particular on a line chart ($j=1$) the parameter degree is at most one.
\end{claim}

\begin{proof}
A maximal minor is multilinear in the rows; exactly the $j$ rows corresponding to a basis of $\Phi$ vary,
and they vary affinely in the chart coordinates.
\end{proof}

\begin{thm}[bit complexity]\label{thm:complexity}
Assume (i) $c=2$, (ii) simple spectrum, (iii) the split condition, and (iv) \textbf{$p$ fixed}, so that
$n-k=\rank[C;CA]\le2p$ is bounded. Then Algorithm~\ref{alg:decide} runs in $\poly(n,\tau)$ bit operations,
with arithmetic in number fields of degree at most $k^2$.
\end{thm}

\begin{proof}
There are $O(k^2)$ planes; for each, the realisable image spaces live in $V_\varphi=\varphi(T_D)$, whose
dimension is at most $2p$ and hence bounded; so the number of charts and the number of parameters per
chart are bounded independently of $k$ (at most $O(p^2)$ parameters, by Remark~\ref{rem:imagesize}). On
such a chart the pencil has $\deg_\lambda\le\omega$ and heights $\poly(n,\tau)$, and the gcd of the
maximal minors over $L_D(\text{chart})[\lambda]$ costs
$\poly(\omega,p,[L_D:\mathbb Q],\text{heights})$ by the standard subresultant bound \cite{vzgg13}.
Summing over the $O(k^2)$ planes gives $\poly(n,\tau)$.
\end{proof}

\begin{remark}[the two boundaries, and why $p$ must be fixed]\label{rem:boundaries}
Hypothesis (iv) of Theorem~\ref{thm:complexity} is not cosmetic. Without a bound on
$n-k=\rank[C;CA]$ the space $V_\varphi=\varphi(T_D)$ grows with $k$, the fibre parameter has up to
$p^2$-many free parameters, and the chartwise gcd is no longer a polynomial-time primitive. Note also
that $c=2$ alone gives $\omega=k-2$, so $D^\circ$ is \emph{large}: it is the bound on $n-k$, not the
smallness of $c$, that controls the fibre. Likewise, without simplicity of the spectrum the base is
positive-dimensional (a Grassmannian of dimension $2(d_\lambda-2)=O(n)$) and the number of charts is no
longer polynomial. With $c$ allowed to vary, the number of components is $\binom{k}{c}$, which is
exponential. The present method therefore gives $\poly(n,\tau)$ only under fixed $p$ with simple
spectrum, and at best $n^{O(cn+p^2)}$ in general.
\end{remark}

\subsection{A worked example}\label{sec:example}

This subsection carries one instance through the reduction and the certificate: the compressed data,
the admissible kernel direction, the image space, the pencil, and the gcd certificate. All entries are
exact integers, and the computation is the one behind the $c=2$ side of Theorem~\ref{thm:nonmono},
with the intermediate matrices displayed.

\emph{The instance.} $n=8$, $p=2$, and
\begin{equation}\label{eq:exA}
\begin{gathered}
A=\begin{pmatrix}
0 & 0 & 1 & 0 & 0 & 0 & 0 & 0\\
0 & 0 & 0 & 1 & 0 & 0 & 0 & 0\\
1 & -2 & 1 & 0 & -1 & 1 & 1 & -1\\
-1 & 2 & 1 & 0 & 2 & -2 & -1 & 2\\
0 & 0 & 0 & 0 & 1 & 0 & 0 & 0\\
0 & 0 & 0 & 0 & 0 & 1 & 0 & 0\\
0 & 0 & 0 & 0 & 0 & 0 & 2 & 0\\
0 & 0 & 0 & 0 & 0 & 0 & 0 & 3
\end{pmatrix},\\
C=\begin{pmatrix}
1&0&0&0&0&0&0&0\\ 0&1&0&0&0&0&0&0
\end{pmatrix}.
\end{gathered}
\end{equation}
Here $\rank[C;CA]=4$, so $k=4$, and
\[
K_0=\ker[C;CA]=\operatorname{span}\{e_5,e_6,e_7,e_8\},\qquad
N=\begin{pmatrix}0_{4\times4}\\ I_4\end{pmatrix}.
\]
The compression is $M=\pi_{K_0}A|_{K_0}=\operatorname{diag}(1,1,2,3)$ and the coupling is
\[
A_{21}=\begin{pmatrix}
0&0&0&0\\ 0&0&0&0\\ -1&1&1&-1\\ 2&-2&-1&2\\ 0&0&0&0\\ 0&0&0&0\\ 0&0&0&0\\ 0&0&0&0
\end{pmatrix}.
\]
The coupling has rank two, but its restriction to the direction below has rank one.

\emph{The direction.} In the coordinates of $K_0$ fixed by $N$, take
$D=\operatorname{span}\{e_1+e_2,\,e_3\}\le K_0$. Then $M(e_1+e_2)=e_1+e_2$ and $Me_3=2e_3$, so
$M(D)\subseteq D$ and $D\in\mathcal I_2(M)$ with $\dim D=c=2$. Since $L_0=0$ the prescribed-row
conditions read $P_{K_0}(V_0)=0$ and $W_0=0$, which hold trivially, so $D\in\mathcal B_{\rm pres}$.

\emph{The fibre.} Here $A_{21}D=\operatorname{span}\{(0,0,1,-1,0,0,0,0)^{\mathsf T}\}$, so
\[
r(D)=1,\qquad \dim T_D=(n-k)-r(D)=4-1=3.
\]
With $t=1$ (i.e.\ $q=\omega+t=2+1=3$), $w:=(0,0,1,1,0,0,0,0)^{\mathsf T}$ and
$W=\operatorname{span}\{w\}\subseteq T_D$, one computes
\[
\varphi_w=(1,-1,0,1)\in D^\circ,\qquad \Phi(W)=\operatorname{span}\{(1,-1,0,1)\},
\]
so $\dim\Phi(W)=1=\min(t,2p)$: the image parameter is a line, as in Definition~\ref{defn:charts}. The
resulting design row space is
\[
\begin{aligned}
V=D^{\perp_{K_0}}\oplus W
&=\operatorname{span}\bigl\{(0,0,0,0,-1,1,0,0),\\
&\qquad (0,0,0,0,0,0,0,1),\\
&\qquad (0,0,1,1,0,0,0,0)\bigr\},
\end{aligned}
\]
of dimension $q=3$; the design is $K_0$-split by construction.

\emph{The pencil.} With these three vectors as the rows of $\mathcal L$,
\[
\begin{gathered}
H=\mathcal L N=\begin{pmatrix}-1&1&0&0\\ 0&0&0&1\\ 0&0&0&0\end{pmatrix},\\
H_A=\mathcal L AN=\begin{pmatrix}-1&1&0&0\\ 0&0&0&3\\ 1&-1&0&1\end{pmatrix}.
\end{gathered}
\]
The invariants are $\omega=\rank H=2$, $\dim(K_0\cap V^\perp)=2=c$ (exact kernel), and
$\row H_A\subseteq\row H$, i.e.\ \eqref{eq:5a} holds. The nonzero maximal $2\times2$ minors of
$\lambda H-H_A$ are
\[
\begin{gathered}
-(\lambda-3)(\lambda-1),\quad (\lambda-3)(\lambda-1),\quad \lambda-1,\\
1-\lambda,\quad \lambda-3,\quad 3-\lambda .
\end{gathered}
\]

\emph{The certificate.} The minors $\lambda-1$ and $\lambda-3$ are coprime, with Bézout identity
\[
\tfrac12\bigl((\lambda-1)-(\lambda-3)\bigr)=1 ,
\]
so the gcd of all maximal minors is $1$ and condition~\eqref{eq:5b} holds. This is a direct
certificate for the displayed design; it does not require the algorithm. Indeed a feasible design of
order $q=3$ with $c=2$ exists, whereas by Theorem~\ref{thm:nonmono} no feasible design of order $3$
with $c=1$ does.

\emph{Scope.} The eigenvalue $1$ of $M$ has a two-dimensional eigenspace, so $\mathcal I_2(M)$ is
positive-dimensional (Remark~\ref{rem:nonsimple}) and the finite enumeration of
Theorem~\ref{thm:base} does not apply. The proofs of Theorems~\ref{thm:correct} and
\ref{thm:complexity} use simple spectrum in the sense of distinct eigenvalues; this instance
therefore illustrates the reduction and the certificate format, not the algorithm's execution or
those proofs. The distinction is worth recording, because it shows that feasibility itself is not
confined to the distinct-eigenvalue case --- only the correctness and cost arguments are.

\section{Discussion and conclusion}\label{sec:conclusion}

This section collects what the paper does not claim and records the evidence behind what it does claim.
Part A consolidates the limitations, Part B the evidence.

\subsection{Limitations}\label{sec:conclusion:limits}

\begin{enumerate}[label=\arabic*.,leftmargin=*,itemsep=1pt,topsep=2pt]
\item All structural results are for the split problem \textbf{S}; the unrestricted problem \textbf{U} is
only reduced to a fibre description (Theorem~\ref{thm:ufibre}) and its invariant-kernel structure is
lost.
\item The decision algorithm and the complexity bound require simple spectrum; the non-simple case is
left open, and the worked example of Section~\ref{sec:example}, whose compression has a repeated
eigenvalue, illustrates the certificate format rather than these two results.
\item No uniformity in $c$ is claimed; the honest general bound is $n^{O(cn+p^2)}$.
\item The general invariant-plane variety is used only through its defining equations; no rational
parametrization is claimed.
\end{enumerate}

The disclaimers inside the body are part of the same statement and are not weakened here:
Remark~\ref{rem:nonsimple} is the disclaimer attached to the classification of
Theorem~\ref{thm:base}, Remark~\ref{rem:boundaries} lists the two boundaries of the complexity bound, and
Remark~\ref{rem:open} lists the two questions on general $c$ that we do not answer: (a) whether the
quantifier over $\Phi$ can be replaced by a rank or separation condition on the induced maps of
$T_D/W_0$, and (b) whether the layer index can be traded for $\dim U$ with $U=\pi_1(W)$, since
$c=k-\dim U$. Nothing is claimed about the complexity of the general decision problem.

\subsection{Reproducibility and evidence}\label{sec:conclusion:evidence}

This paper has \emph{no} machine-checked formalisation. Unlike companion manuscripts that carry Lean
developments, no statement of Sections~\ref{sec:layers}--\ref{sec:decision} is formalised at the time of
writing, and nothing in those sections should be read as machine-verified.

The evidence behind the concrete assertions is of two kinds. First, the proofs of
Sections~\ref{sec:layers}--\ref{sec:decision} are paper-level and complete, with one exception:
the height chain of Claim~\ref{clm:arith} is a sketch in the body, made explicit step by step in
Appendix~\ref{app:height}; the constants there are standard and deliberately not optimized, so no
single numerical constant is claimed. Second, the three places where the paper asserts something
about a specific system --- the instance of Theorem~\ref{thm:nonmono}, the worked example of
Section~\ref{sec:example}, and the classification of Table~\ref{tab:planes} --- rest on exact
rational arithmetic, not on floating-point computation. In the proof of Theorem~\ref{thm:nonmono}
the nonzero maximal $2\times2$ minors are expanded over $\mathbb Q$ and their gcd is $1$ precisely
because the displayed combination $\tfrac12\bigl((\lambda-1)-(\lambda-3)\bigr)$ is a rational
combination of two of them, and on the $c=1$ side the gcd $(\lambda-1)(\lambda-2)(\lambda-3)$ is
read off the displayed integer matrix $H$ for symbolic $s$; both computations involve no algebraic
extension and no rounding. In the worked example of Section~\ref{sec:example} the data $A,C$ are
integral, and $N$, $M$, $A_{21}$, the design basis, $H$, $H_A$, all $2\times2$ minors of the pencil
and their gcd are computed over $\mathbb Q$; six minors are nonzero, the displayed Bézout identity
is exact, and again no algebraic extension or rounding occurs, so the certificate is reproducible in
any exact arithmetic system. Table~\ref{tab:planes} is a classification rather than a computation:
its two counts are the combinatorial expressions of Theorem~\ref{thm:base}, and its third column is
the field statement of Claim~\ref{clm:arith}.

\appendices
\small
\section{Proof of the height propagation}\label{app:height}

This appendix records the chain used in Claim~\ref{clm:arith} with the standard bounds made
explicit. Throughout, $h(\cdot)$ is the maximum bit length of the entries of a matrix or of the
coefficients of a polynomial.

\emph{Step 0 (input).} $h(A),h(C)\le\tau$.

\emph{Step 1 ($CA$ and $[C;CA]$).} Entrywise, $(CA)_{ij}$ is a sum of $n$ products, so
$h(CA)\le 2\tau+\lceil\log_2 n\rceil$. Hadamard's bound gives, for every square submatrix $S$ of
$[C;CA]$ of order $d$, $|\det S|\le(\sqrt{n}\,2^{2\tau+\log n})^{d}$, i.e.\
$h(\det S)\le d\,(2\tau+\log n+\tfrac12\log n)$.

\emph{Step 2 (kernel basis $N$).} A basis of $\ker[C;CA]$ can be chosen so that each entry is a
cofactor of $[C;CA]$, hence of order at most $2p$; by Step 1,
\[
h(N)\le 2p\,(2\tau+\log n+\tfrac12\log n).
\]

\emph{Step 3 (Gram matrix and its inverse).} $G=N^{\mathsf T}N$ has entries that are sums of $n$
products, so $h(G)\le 2h(N)+\log n$. Since $G$ is a positive definite integer matrix, $\det G\ge1$,
and the entries of $G^{-1}$ are cofactors over $\det G$, so
\[
h(G^{-1})\le k\,(2h(N)+\log n+\tfrac12\log k).
\]

\emph{Step 4 (compression and coupling).} $M=(N^{\mathsf T}N)^{-1}N^{\mathsf T}AN$ and
$A_{21}=P_{K_0^\perp}AN$ are obtained from $G^{-1}$, $N$ and $A$ by polynomially many ring
operations, so
\[
h(M),\,h(A_{21})\le 4h(N)+2h(G^{-1})+\tau+O(\log n),
\]
which is $\poly(n,\tau)$.

\emph{Step 5 (characteristic and minimal polynomial).} The coefficient of $t^{\,k-i}$ in $\chi_M$ is
a sum of principal $i\times i$ minors, hence at most $\binom{k}{i}(\sqrt{k}\,2^{h(M)})^{i}$ in
absolute value; so $h(\chi_M)\le k\,h(M)+O(k\log k)$. The minimal polynomial divides $\chi_M$.

\emph{Step 6 (irreducible factors).} By Mignotte's bound, a divisor $f$ of a degree-$k$ integer
polynomial $g$ satisfies $\|f\|_\infty\le\binom{k}{\lfloor k/2\rfloor}\|g\|_\infty$, so
\[
h(f)\le h(\chi_M)+k .
\]

\emph{Step 7 (roots).} Represent a root $\alpha$ of $f$ by the class of $t$ in
$\mathbb Q[t]/(f)$; its height is bounded in terms of $h(f)$. Cauchy's bound gives
$|\alpha|\le2\|f\|_\infty^{1/\deg f}$, and the logarithmic height satisfies
$h(\alpha)\le h(f)+\log\deg f$.

\emph{Step 8 (eigenvectors).} Solving $(M-\alpha I)v=0$ over $L=\mathbb Q(\alpha)$ by Cramer
expresses each coordinate of $v$ as a quotient of two minors of order $k-1$ of $M-\alpha I$. Those
minors are polynomials in $\alpha$ of degree at most $k-1$ whose coefficients involve $k!$ products
of entries of $M$; hence
\[
h(v)\le (k-1)\bigl(h(M)+h(\alpha)\bigr)+O(k\log k).
\]
The denominators are nonzero because the eigenvalues are distinct.

\emph{Step 9 (subspaces and the image map).} The bases of $D$, $D^{\perp_{K_0}}$ and $T_D$ are
kernels and orthogonal complements of matrices with entries of the heights above, so they inherit an
$O(k)$ factor. The functionals $\varphi$, the prescribed image $\Phi_0$ and the entries of the cleared
pencil are obtained from these by polynomially many field operations, so their heights are again
$\poly(n,\tau)$.

\emph{Step 10 (gcd).} The maximal minors have $\lambda$-degree at most $\omega\le k$ and heights at
most $H=\poly(n,\tau)$. Each subresultant is a determinant of order at most $2k$ whose entries are
coefficients of the maximal minors, so Hadamard's bound gives bit length $O(k(H+\log k))$, which is
$\poly(n,\tau)$; the same bound applies to the gcd (see \cite{vzgg13}).

Combining Steps 0--10 yields Claim~\ref{clm:arith}. The constants above are not optimized: each step
multiplies the height by a factor polynomial in $k$ and in the degree of the field, and the field
degree is at most $k^2$; no attempt has been made to minimize the exponents.

\normalsize

\begin{thebibliography}{99}
\bibitem{companion1} Deciding functional observer feasibility: compression, certificates, and
parameterized algorithms. Companion manuscript, 2026.
\bibitem{companion2} Minimum order and minimum-order construction for functional observers. Companion
manuscript, 2026.
\bibitem{companion3} Functional-observer feasibility in the hyperplane layer. Companion manuscript, 2026.
\bibitem{fer26a} A. Fernando. \emph{Functional observer design: existence and order spectrum}.
arXiv:2609.33615, 2026 (preprint).
\bibitem{fer26b} A. Fernando. \emph{Functional observers: a chartwise characterisation}. arXiv:2609.37992,
2026 (preprint).
\bibitem{bruns88} W. Bruns, U. Vetter. \emph{Determinantal Rings}. Lecture Notes in Mathematics 1327.
Berlin, Germany: Springer, 1988.
\bibitem{fulton97} W. Fulton. \emph{Young Tableaux}. Cambridge, U.K.: Cambridge University Press, 1997
(degeneracy loci, Ch.~14).
\bibitem{bpr06} S. Basu, R. Pollack, M.-F. Roy. \emph{Algorithms in Real Algebraic Geometry}, 2nd ed.
Berlin, Germany: Springer, 2006.
\bibitem{vzgg13} J. von zur Gathen, J. Gerhard. \emph{Modern Computer Algebra}, 3rd ed. Cambridge, U.K.:
Cambridge University Press, 2013 (subresultants, gcd degree bounds).
\bibitem{emre80} E. Emre, M. L. J. Hautus. A polynomial characterization of $(A,B)$-invariant and
reachability subspaces. \emph{SIAM J. Control Optim.}, vol.~18, no.~4, 1980.
\end{thebibliography}

\end{document}
